\documentclass[10pt,a4paper]{amsart}
\usepackage[margin=19mm]{geometry}
\usepackage{amsmath,amssymb,mathtools}
\usepackage[scr=rsfs,cal=euler]{mathalfa}
\usepackage{xfrac}
\usepackage[unicode,hidelinks,bookmarks=false]{hyperref}
\newcommand{\ie}{i.e.\ }
\newcommand{\cf}{cf.\ }
\newcommand{\resp}{respectively\ }
\newcommand{\Subsection}{\subsection}
\providecommand{\nobreakdash}{\nobreak\hbox{-}}
\setlength{\emergencystretch}{2em}
\multlinegap0pt
\usepackage{bm}
\usepackage{bbm}
\usepackage{dsfont}
\usepackage{xspace}
\usepackage{cancel}
\usepackage{mathtools}
\usepackage{amsthm}
\makeatletter
\@ifundefined{theorem}{\newtheorem{theorem}{Theorem}}{}
\@ifundefined{lemma}{\newtheorem{lemma}[theorem]{Lemma}}{}
\makeatother
\newcommand{\R}{{\mathbb R}}
\newcommand{\Z}{{\mathbb Z}}
\title[The Bombieri-Pila determinant method]{The Bombieri-Pila determinant method}
\author{Thomas F. Bloom, Jared Duker Lichtman}
\date{Source: arXiv:2312.12890v3 (CC BY 4.0)}
\begin{document}
\maketitle

\noindent\emph{Source and licence.} The Bombieri-Pila determinant method. Version arXiv:2312.12890v3 (CC BY 4.0), distributed under the Creative Commons Attribution 4.0 International licence, as recorded in the arXiv licence metadata for this exact version and shown on the abstract page https://arxiv.org/abs/2312.12890v3. Full rights notice: licenses/arxiv-2312.12890-CC-BY-4.0.txt.\par\smallskip

\noindent\textit{Editorial scope.} Counted unit: the Theorem whose source label is \texttt{thm:GNbound} in the article (source0.tex lines 516--521, source label \texttt{thm:GNbound}), delivered together with the article's own complete original proof of it (source0.tex lines 522--560, one \texttt{proof} environment of 39 lines). No mathematics is written, repaired or re-derived: statement and proof are the article's own text line for line. Required context retained uncounted: lem-intbound1 (lines 374--381); eq:Dcompute (lines 394--396); lem:6 (lines 479--485); lem:7 (lines 492--500). Every definition, display and label of those lines is retained unchanged. Omitted from this package and not redistributed: 1--373, 382--393, 397--478, 486--491, 501--515, 561--741. Nothing inside a retained excerpt is omitted or altered: every retained body is delivered exactly as the article's lines are written, so no omission receipt of any kind accompanies this package. Citations: the retained excerpts carry no citation command at all, so no bibliography entry of the article is transcribed and no reference is redistributed. Reference adaptation: none was needed and none was made; every reference inside the delivered unit resolves inside it, so no unresolved reference or citation is produced. Printed slips kept literal: no printed word, symbol, hypothesis or number of the counted statement or of its proof was altered. Layout and class: the article's own class is \texttt{amsart} with packages including fullpage, amsfonts, amsmath, amsthm, amssymb, enumitem, comment, physics, mathtools; the wrapper is the lane's pinned \texttt{amsart} editorial wrapper at 10pt on A4, which restates the theorem-like environments the retained text needs and adds the article's own macro definitions that the retained text needs (\texttt{\textbackslash R}, \texttt{\textbackslash Z}), with extra wrapper packages (bm, bbm, dsfont, xspace, cancel, mathtools). The delivered numbering is the unit's own. Assets: no retained span contains a figure environment, a graphics-inclusion command, a PDF-page inclusion, a table or a TikZ picture; no image file is copied into this package, the package declares no asset dependency and no image file is redistributed with it. Licence: version arXiv:2312.12890v3 of this article is distributed under the Creative Commons Attribution 4.0 International licence, as recorded in the arXiv licence metadata for this exact version and shown on the abstract page https://arxiv.org/abs/2312.12890v3. A full rights notice is delivered at licenses/arxiv-2312.12890-CC-BY-4.0.txt. Characters: every retained excerpt is pure ASCII, so the delivered engine, XeLaTeX with its default Latin Modern text font, reports no missing character.
\begin{lemma}\label{lem-intbound1}
Let $\ell \geq d\geq 2$. Let $I\subseteq [0,N]$ be a closed interval, and $f\in C^{\infty}(I)$ with $F(x,f)=0$ for some irreducible polynomial $F\in \R[x,y]$ of degree $d\geq 2$. Suppose that $N>0$ and $\delta\geq 1/\lvert I\rvert$ satisfy
\[\left\lvert \frac{f^{(i)}(x)}{i!}\right\rvert\leq N\delta^i\qquad {\rm for \ all }\quad x\in I,\]
for every $0\leq i<D=d(\ell-d+1)$. Then we have
\begin{align*}
\left\lvert\big\{(x,f(x)) : x\in I\big\}\cap \Z^2\right\rvert \ll(d\ell)^2 N^{\frac{1}{d}+O(\frac{1}{\ell})} \delta\lvert I\rvert.
\end{align*}
\end{lemma}

\begin{equation}
D =\lvert \mathcal{M}\rvert = \sum_{d\le h\le \ell} d  =d(\ell-d+1) = d\ell+O(d^2) \label{eq:Dcompute}
\end{equation}

\begin{lemma}\label{lem:6}
Suppose $F(x,y)\in\R[x,y]$ is an irreducible polynomial of degree $d\ge2$. Let $I$ be an interval and $f(x)\in C^\infty(I)$ satisfy $F(x,f)=0$. Let $A_1,\ldots,A_k>0$. We may partition $I$ into $O(k^2d^2)$ many subintervals $I_\nu$ such that, for each interval $I_\nu$ and each $1\leq i\leq k$,
\begin{align*}
{\rm (i)}&\qquad |f^{(i)}(x)| \le A_i \qquad {\rm for \ all }\quad x\in I_\nu, \qquad {\rm or}\\
{\rm (ii)}&\qquad |f^{(i)}(x)| \ge A_i \qquad {\rm for \ all }\quad x\in I_\nu.
\end{align*}
\end{lemma}

\begin{lemma}\label{lem:7}
Let $k\geq 1$, $0<\delta<1$, and $X>0$. Let $I$ be an interval and $f\in C^k(I)$. If
\begin{align*}
\Big|\frac{f^{(i)}(x)}{i!}\Big| \le X\delta^i
\quad\textrm{ for all }0\leq i<k
\quad\textrm{ while }\quad\Big|\frac{f^{(k)}(x)}{k!}\Big| \ge X\delta^k
\end{align*}
for every $x\in I$, then $\lvert I\rvert \leq 2/\delta$.
\end{lemma}

\begin{theorem}\label{thm:GNbound}
Let $d\geq 2$ and $N\geq 1$ be some integer sufficiently large in terms of $d$, and let $I$ be some interval of length $N$. Let $f\in C^\infty(I)$ be such that $\lvert f'(x)\rvert \leq 1$ for all $x\in I$. If there exists some irreducible polynomial $F(x,y)\in \R[x,y]$ of degree $d\geq 2$ such that $F(x,f(x))= 0$, identically as a function of $x$, then 
\begin{align*}
\left\lvert\big\{ (x,f(x)) : x\in I\big\}\cap \Z^2\right\rvert \le (\log N)^{O(d)}N^{\frac{1}{d}}.
\end{align*}
\end{theorem}

\begin{proof}
For any real $N>0$ define
\[
G(N):= \sup_{|I|\le N} \sup_{\substack{f\in C^\infty(I)\\ \lvert f'\rvert\leq 1}}\left\lvert \big\{(x,f(x)) : x \in I \big\}\cap \Z^2 \right\rvert.\]
We shall prove that for any integer $\ell \geq d$ there exists some $K=K(d,\ell)$ satisfying $2\leq K\leq \ell^{O(1)}$ such that for any $N>0$ we have
\begin{equation}\label{eq-recur}
G(N) \leq K^{O(d)} N^{\frac{1}{d}+O(\frac{1}{\ell})} + KG(K^{-2d}N).
\end{equation}
We first show how \eqref{eq-recur} implies Theorem \ref{thm:GNbound}. Fix some large $N>0$ and $\ell\geq d$ (which will depend on $N$). Iteratively applying \eqref{eq-recur}, since $K(K^{-2d})^{\frac{1}{d}}= K^{-1}\leq  1/2$ (assuming $\ell$ is large enough), it follows by induction that for any $n\geq 1$
\[G(N) \leq K^{O(d)}N^{\frac{1}{d}+O(\frac{1}{\ell})}\sum_{0\leq j<n}2^{-j}+K^nG(K^{-2nd}N).\]
In particular, if we choose $n$ large enough such that $K^{2nd}\in [N,K^{2d}N)$, we have $K^n \leq KN^{1/2d}$ and $G(K^{-2nd}N)\leq G(1)\le 1$, so that
\begin{align*}
G(N) & \ll K^{O(d)}N^{\frac{1}{d}+O(\frac{1}{\ell})}+KN^{1/2d} \ll \ell^{O(d)}N^{\frac{1}{d}+O(\frac{1}{\ell})} \ll(\log N)^{O(d)}N^{\frac{1}{d}},
\end{align*}
recalling $K\leq \ell^{O(1)}$ and choosing $\ell=\lceil \log N\rceil$. Thus \eqref{eq-recur} implies the result.

Now to prove \eqref{eq-recur} itself, we fix some interval $I$ of length $N$ and $f\in C^\infty(I)$ with $\lvert f'\rvert\leq 1$. Without loss of generality, translating the graph of $f$ by an integer if necessary, we may assume that $I=[0,N]$ and $\lvert f(x)\rvert \leq N$ for all $x\in I$. Fix some $\ell\geq d$ and let $\delta=\delta(d,\ell)\in(2/N,1)$ be some quantity to be chosen later. 

We shall apply Lemma \ref{lem:6} with $A_i=N\delta^i$, and  recall $D=|\mathcal M|\ll d\ell$ from \eqref{eq:Dcompute}. Indeed, by Lemma \ref{lem:6} we may partition $I$ into at most $O(d^2D^2)$ subintervals $I_\nu$ such that, for each $I_\nu$ and each $1\le i<D$, either 
\begin{align*}
{\rm (i)} & &\Big|\frac{f^{(i)}(x)}{i!}\Big| \ &\le N\delta^i \qquad {\rm for \ all }\quad x\in I_\nu,\textrm{ or }\\
{\rm (ii)} & &\Big|\frac{f^{(i)}(x)}{i!}\Big| \ &\ge N\delta^i  \qquad {\rm for \ all }\quad x\in I_\nu.
\end{align*}
Suppose first that $I_\nu$ satisfies $\lvert I_\nu\rvert \geq 1/\delta$ and (i) holds for every $1\leq i<D$. Then applying Lemma~\ref{lem-intbound1} we have
\begin{align*}
\left\lvert \big\{ (x,f(x)) : x\in I_\nu\big\}\cap \Z^2\right\rvert\ll
(d\ell)^2 N^{\frac{1}{d}+O(\frac{1}{\ell})}\delta|I_\nu|.
\end{align*}

Otherwise, either $\lvert I_\nu\rvert \leq 1/\delta$, or there exists a (minimal) $1\le k<D$ such that (ii) holds for $k$, but (i) holds for all $i<k$. In this case, Lemma \ref{lem:7} implies $|I_\nu| \le 2\delta^{-1}$.

Summing the contribution from each $I_\nu$ we deduce that 
\begin{align*}
\left\lvert \big\{(x,f(x)) : x \in I \big\}\cap \Z^2 \right\rvert & =  \sum_{\nu\ll d^2D^2}\left\lvert \big\{(x,f(x)) : x \in I_\nu \big\}\cap \Z^2 \right\rvert\\
& \ll (d\ell)^2 N^{\frac{1}{d}+O(\frac{1}{\ell})}\delta \sum_{\nu\ll d^2D^2}|I_\nu|+ d^2D^2G(2\delta^{-1}) \\
&\ll d^4\ell^2 \left( N^{\frac{1}{d}+O(\frac{1}{\ell})}\delta N +G(2\delta^{-1})\right),
\end{align*}
using the fact that $\sum_\nu \lvert I_\nu\rvert=\lvert I\rvert\leq N$. This gives \eqref{eq-recur} for some $K\ll d^4\ell^2\leq\ell^{O(1)}$, after choosing $\delta = 2K^{2d}/N$.
\end{proof}

\end{document}
